\section{The deflating spacetime}
\label{sec:geometry}

This section derives the geometry that the hypothesis implies: the curvature, the field equations
that the metric does and does not solve, its symmetries, and the motion of a classical particle.
Every formula of this section was obtained twice by exact computer algebra, by the program of this
paper and by an independent script (Appendix~\ref{app:checks}); the steps are short enough to be
repeated by hand, and the text says how.

\subsection{Metric, frame and volume}

Throughout, the metric is written in the sign convention of the hypothesis (\ref{eq:hyp}):
\emph{space-like components are positive and time-like components are negative}. The flat metric is
\begin{equation}
\etam=\mathrm{diag}\SIG{(+1,+1,+1,+1,-1,-1,-1,-1)}{(+1,+1,+1,+1,-1,+1,+1,+1)},
\label{eq:etam}
\end{equation}
with the coordinates in the order $x0,\dots,x7$: \SIG{four space-like coordinates $x0,\dots,x3$ and
four time-like ones, $x4,\dots,x7$.}{seven space-like coordinates and the one time-like coordinate
$x4$.} Define three \emph{scale factors}, one for each deflating coordinate,
\begin{equation}
b_t(x4)=\ct\,e^{-\rt\,x4},\qquad t=5,6,7,
\label{eq:scale}
\end{equation}
with $\mathrm{rate}_5=\gff$, $\mathrm{rate}_6=\gss$, $\mathrm{rate}_7=\gvv$. The metric of the
deflating spacetime is
\begin{equation}
\gm_{\mu\nu}=\mathrm{diag}\big(+1,+1,+1,+1,-1,\,\SIG{-}{+}b_5^2,\,\SIG{-}{+}b_6^2,\,\SIG{-}{+}b_7^2\big),
\label{eq:metric}
\end{equation}
whose last three components are those of Eq.~(\ref{eq:hyp}). The superscript \SIG{$(4+4)$}{$(7+1)$}
marks the metric of this spacetime in this sign convention; a component of the metric always
carries it, so that $\gm_{55}$ cannot be taken for the rate $\gff$. Greek indices label the
coordinates. A
coordinate interval $dxt$ has the proper extent $b_t\,|dxt|$, which shrinks by the factor $e$ in the
time $1/\rt$. \SIG{Because $x5,x6,x7$ are time-like, what shrinks is a duration measured along an
additional time.}{Because $x5,x6,x7$ are space-like, what shrinks is a length.}

Capital indices $A,B$ label an orthonormal frame, eight vector fields $e_A=e_A{}^{\mu}\partial_\mu$
with $\gm_{\mu\nu}e_A{}^{\mu}e_B{}^{\nu}=\etam_{AB}$. We take the frame along the coordinates,
$e_A{}^{\mu}=\mathrm{diag}(1,1,1,1,1,b_5^{-1},b_6^{-1},b_7^{-1})$. The volume element is $\Vd\,d^8x$
with
\begin{equation}
\Vd(x4)=b_5b_6b_7\propto e^{-\gsum\,x4},\quad\gsum=\gff+\gss+\gvv .
\label{eq:Vd}
\end{equation}
$\Vd$ is the proper volume of a unit coordinate cell of the deflating coordinates, and the
\emph{total rate} $\gsum$ is the rate at which it shrinks. Most results below depend on the three
rates only through $\gsum$.

\subsection{Curvature}
\label{sec:curvature}

The only components of the metric that depend on a coordinate are the three $\gm_{tt}$, and they
depend on $x4$ alone. The Christoffel symbols
$\Gamma^{\rho}{}_{\mu\nu}=\tfrac12g^{\rho\sigma}(\partial_\mu g_{\sigma\nu}+\partial_\nu g_{\sigma\mu}-\partial_\sigma g_{\mu\nu})$
that do not vanish are therefore
\begin{equation}
\Gamma^{t}{}_{4t}=-\rt,\qquad\Gamma^{4}{}_{tt}=\SIG{+}{-}\rt\,b_t^{2},
\end{equation}
without summation over $t$. The Riemann tensor has, in the orthonormal frame, the components
\begin{equation}
R^{t}{}_{4t4}=-\rt^{2},\quad R^{t}{}_{sts}=\SIG{-}{+}\rt\,\mathrm{rate}_s\ \ (t\ne s),
\label{eq:riemann}
\end{equation}
and those related to them by symmetry. No component carries an index of $x0,\dots,x3$: these four
coordinates span a flat space, and the curvature lives in the four coordinates $x4,\dots,x7$. With
$\gsq=\gff^2+\gss^2+\gvv^2$ and $\gprod=\gff\gss+\gff\gvv+\gss\gvv$, the curvature scalar is the
constant $R=\gsq+\gsum^{2}$, and the Einstein tensor
$\cG^{\mu}{}_{\nu}=R^{\mu}{}_{\nu}-\tfrac12R\,\delta^{\mu}_{\nu}$ is diagonal with
\begin{equation}
\begin{aligned}
\cG^{4}{}_{4}&=-\gprod,\qquad\cG^{j}{}_{j}=-\tfrac12(\gsq+\gsum^{2}),\\
\cG^{t}{}_{t}&=\rt\,\gsum-\tfrac12(\gsq+\gsum^{2}),
\end{aligned}
\label{eq:einstein}
\end{equation}
with $j=0,1,2,3$. In the sign convention used here these mixed components are the same on the $4+4$
and on the $7+1$ spacetime. (They follow from the sectional curvatures of Appendix~\ref{app:curvature}:
$\cG^{A}{}_{A}$ is minus the sum of the sectional curvatures of the planes that do not contain the
direction $A$.)
The Einstein tensor does not vanish, and it is not a multiple of the unit matrix unless all rates
vanish. \emph{The deflating spacetime is not a solution of Einstein's equations without matter, with
or without a cosmological constant.}

\subsection{A solution of Einstein-Gauss-Bonnet gravity}
\label{sec:egb}

In more than four dimensions the general field equations of second order for the metric are those of
Lovelock \cite{Lovelock1971}. Their first term beyond Einstein's is the Lanczos tensor
$\cH^{\mu}{}_{\nu}$ of the Gauss-Bonnet term \cite{BoulwareDeser1985}, and the equations without
matter read
\begin{equation}
\cG^{\mu}{}_{\nu}+\Lam\,\delta^{\mu}_{\nu}+\aGB\,\cH^{\mu}{}_{\nu}=0,
\label{eq:egb}
\end{equation}
with a cosmological constant $\Lam$ and a coupling $\aGB$. For the metric
(\ref{eq:metric}) we find
\begin{equation}
\begin{aligned}
\cH^{j}{}_{j}&=-4\,\gff\,\gss\,\gvv\,\gsum\quad(j=0,1,2,3),\\
\cH^{4}{}_{4}&=\cH^{t}{}_{t}=0,
\end{aligned}
\label{eq:lanczos}
\end{equation}
and the Lovelock tensor of the third order, the last that can differ from zero in eight dimensions,
vanishes identically. The components of Eq.~(\ref{eq:egb}) along $x4,\dots,x7$ therefore contain no
Gauss-Bonnet term and require the four corresponding components of the Einstein tensor to be equal,
which holds if and only if the three rates are equal. The components along $x0,\dots,x3$ then fix the
coupling. The result:
\begin{equation}
\gff=\gss=\gvv:\qquad\Lam=3\,\gff^{2},\qquad\aGB=-\frac{1}{4\,\gff^{2}} .
\label{eq:egbsol}
\end{equation}
\emph{With three equal rates the deflating spacetime is an exact solution of Einstein-Gauss-Bonnet
gravity without matter}, in a theory whose two constants satisfy $4\aGB\Lam=-3$; the rate is then
$\gff=\sqrt{\Lam/3}$. With unequal rates it is not a solution for any $\Lam$ and $\aGB$. Solutions of
this theory in which every scale factor is an exponential of the time have been classified in five
and six dimensions \cite{ChirkovPavluchenkoToporensky2014}; the present one, with four static and
three contracting directions in eight dimensions, is of that kind. In the notation of the
Einstein-Lovelock equations of Sec.~\ref{sec:gravity}, Eq.~(\ref{eq:egbsol}) reads
$\bar\Lambda=3\,\gff^{2}$, $\alpha_2=1/(4\,\gff^{2})$. Section~\ref{sec:gravity} examines the
stability of this solution, with and without matter: it is unstable.

\subsection{Symmetries}

A Killing vector is a vector field along which the metric does not change. Flat spacetime in eight
dimensions has 36. The deflating spacetime has, for any rates, 14: the seven translations along all
coordinates but $x4$; the six rotations of $x0,x1,x2,x3$; and the scaling
\begin{equation}
\cX_{\mathrm{scale}}=\partial_4+\gff\,x5\,\partial_5+\gss\,x6\,\partial_6+\gvv\,x7\,\partial_7,
\label{eq:scaling}
\end{equation}
which shifts the time and stretches the deflating coordinates so as to undo the shrinking. Two equal
rates add the rotation of the two coordinates concerned. Three equal rates add the three rotations of
$x5,x6,x7$ and three further vectors, 20 in all; the four coordinates $x4,\dots,x7$ then form a space
of constant curvature of radius $1/\gff$:
\SIG{minus a four-dimensional hyperbolic space.}{de Sitter space in the half of its flat slicing in
which space contracts.} These lists are complete. (A Killing vector that vanishes at a point is a
Lorentz transformation of the tangent space that leaves the Riemann tensor (\ref{eq:riemann})
unchanged; the dimension of the space of such transformations is 6, 7 or 12, which bounds the number
of Killing vectors by 14, 15 or 20.)

Three consequences are used below. The spacetime is homogeneous: no place and no moment is
distinguished. The constants $\ct$ are not observable by themselves, since a rescaling of $xt$
changes them; the deflation is observable only through an object that fixes a coordinate scale along
$xt$, such as a wave of given coordinate wavelength. And there is no conserved energy, because a
translation in time alone is not a symmetry.

Of the 28 rotations and boosts of flat spacetime, only the six rotations of $x0,\dots,x3$ survive for
general rates (Fig.~\ref{fig:isometries}). Nothing that involves the time survives, in particular
not the boosts of ordinary special relativity: the deflation defines a preferred state of rest, as
the expansion of the universe does.

\begin{figure*}
\includegraphics[width=0.86\textwidth]{fig_isometries}
\caption{\label{fig:isometries}Symmetries of the deflating spacetime. (a) The 28 rotations and boosts
of flat spacetime, one square for each pair of coordinates: those that remain symmetries for any
rates, those that remain symmetries only if the two coordinates deflate equally fast, and those that
are never symmetries. Circles mark the six pairs whose rotations act as internal symmetries on fields
that do not depend on $x0,x5,x6,x7$ (Sec.~\ref{sec:higgs}). (b) The number of Killing vectors.}
\end{figure*}

\subsection{Geodesics}
\label{sec:geodesics}

The translations give seven conserved momenta $P_\mu$ ($\mu\ne4$) of a classical particle. Let
$\Eperp=(m^{2}+\sum_jP_j^{2})^{1/2}$, with the sum over $j=0,1,2,3$, be the energy that a particle of
mass $m$ would have without momentum along the deflating coordinates. The conservation of the norm of the momentum gives the
velocity along the time, with an affine parameter $\varsigma$ along the path:
\begin{equation}
\Big(\frac{dx4}{d\varsigma}\Big)^{2}=\Eperp^{2}\SIG{-}{+}\sum_{t}\frac{P_t^{2}}{b_t(x4)^{2}} .
\label{eq:geodesic}
\end{equation}
The last term grows as $e^{2\,\rt\,x4}$: the momentum $P_t$ belongs to a fixed coordinate wavelength,
and the proper wavelength shrinks.
%%<1644>
With the minus sign the right-hand side vanishes at a \emph{turning value} of the time; for momentum
along $x5$ alone,
\begin{equation}
x4_{\mathrm{turn}}=\frac{1}{\gff}\ln\frac{\Eperp\,\cff}{|P_5|} .
\label{eq:turn}
\end{equation}
The path reaches this value and returns to smaller $x4$ [Fig.~\ref{fig:geodesics}(b)]. A classical
particle with momentum along a deflating time cannot pass the turning value of the ordinary time. Section~\ref{sec:modes} shows the
quantum form of this result. (On the $7+1$ spacetime the sign in Eq.~(\ref{eq:geodesic}) is the
opposite, and the energy of such a particle grows without bound.)
%%</1644>
%%<1671>
With the plus sign the energy of a particle that moves along a deflating coordinate grows without
bound, the reverse of the redshift in an expanding universe. For momentum along $x5$ alone the proper
time that elapses between $x4$ and $x4=\infty$ is
\begin{equation}
\frac{m}{\gff\,\Eperp}\,\mathrm{arsinh}\Big(\frac{\Eperp\,\cff\,e^{-\gff\,x4}}{|P_5|}\Big),
\label{eq:incomplete}
\end{equation}
which is finite [Fig.~\ref{fig:geodesics}(a)]: the coordinates of Eq.~(\ref{eq:metric}) do not cover
the whole future of such a particle. For equal rates this is the known incompleteness of the flat slicing of de Sitter space.
Section~\ref{sec:modes} shows the quantum form of the blueshift. (On the $4+4$ spacetime the sign in
Eq.~(\ref{eq:geodesic}) is the opposite, and such a particle is turned back.)
%%</1671>

\begin{figure*}
\includegraphics[width=0.8\textwidth]{fig_geodesics}
\caption{\label{fig:geodesics}A classical particle of the mass $m=1$ with the momenta $P_1=0.5$ along
$x1$ and $P_5=0.3$ along the deflating coordinate $x5$, for $\gff=1$ and $\cff=0.8$, started at
$x4=-2$. (a) On the $7+1$ spacetime the time $x4$ grows without bound within the finite proper time
3.386. (b) On the $4+4$ spacetime the path reaches the turning value $x4=1.0924$ and returns to
smaller $x4$.}
\end{figure*}

